\documentclass[11pt]{article}
\usepackage[a4paper,margin=28mm]{geometry}
\usepackage{amsmath,amssymb,amsthm}
\usepackage[hidelinks]{hyperref}
\newtheorem{theorem}{Theorem}
\title{A converse for the Edge Clique Cover bound under \(\mathsf{NP}\nsubseteq\mathsf{QP}\)}
\author{Atlas\thanks{The proof in this note was found using ChatGPT.}}
\date{18 November 2026}
\begin{document}
\maketitle

An instance of \textsc{Edge Clique Cover} consists of a finite simple
undirected graph \(G\) with \(n\) vertices and a nonnegative integer \(k\)
given in binary. Vertices are numbered consecutively. The question is whether at most \(k\) cliques cover all
edges of \(G\). The graph is given explicitly. Let \(\mathsf{QP}\) denote
the class of languages decidable in deterministic time
\(2^{O((\log_2(L+2))^c)}\) for some fixed \(c\geq 1\), where \(L\) is the
input length in bits.

\hypertarget{ecc-np-qp-converse}{}
\begin{theorem}\label{thm:ecc-np-qp-converse}
If \textsc{Edge Clique Cover} has no deterministic algorithm running in
time \(2^{k^{O(1)}}n^{O(1)}\), then \(\mathsf{NP}\nsubseteq\mathsf{QP}\).
\end{theorem}

\begin{proof}
Suppose \(\mathsf{NP}\subseteq\mathsf{QP}\). We use a kernel to bound the
length of the input to a quasipolynomial-time algorithm in terms of \(k\).

Every edge can be covered by the clique consisting of its two endpoints,
so we may immediately accept if \(k\geq |E(G)|\). Reading and comparing
the binary threshold takes \(O(\log(k+2))\) time, which the desired bound
absorbs. In the remaining cases the threshold has \(O(\log(n+2))\) bits.

We use the kernel argument of Gramm et al.~\cite[Section 2,
Rules 1 and \(3'\), Theorem 1]{gramm}. To maintain ordinary graph
instances, apply the following rules. Delete isolated vertices.
Delete any connected component consisting of a single edge and
decrease the budget by one, since that component requires one
clique. Reject if the budget becomes negative.

For adjacent vertices \(u,v\) with the same closed neighborhood
outside such a component, delete \(u\). A cover before deletion
restricts to a cover afterward. Conversely, \(v\) has another
neighbor \(x\), so some reduced covering clique covers \(vx\).
Add \(u\) to every clique containing \(v\). These sets are cliques
because the closed neighborhoods agree. They cover \(uv\) and
all edges from \(u\) to their common neighbors.

Repeat the rules until none applies. Each step deletes a vertex
and can be performed in polynomial time. Let \(k'\leq k\) be the
remaining budget. In a cover by at most \(k'\) cliques, each
remaining vertex has a nonempty signature indicating which cliques
contain it. Two vertices with identical signatures would be
adjacent and have identical closed neighborhoods. Thus a reduced
yes-instance has at most \(2^{k'}\) vertices. Reject if there are
more. Otherwise number the remaining vertices consecutively and
encode the graph by an adjacency matrix. The resulting length \(L\), including the budget, is
\(2^{O(k+1)}\).

The language of these encoded instances belongs to \(\mathsf{NP}\).
A certificate can list at most \(\min(k,|E(G)|)\) cliques, since a cover
can be restricted to one chosen clique per edge. A verifier checks that
each listed vertex set is a clique and that every edge belongs to one
of these cliques. The certificate length and verification time are
polynomial in the input length.

By the assumed containment, there is a deterministic algorithm for
this language running in time \(2^{O((\log_2(L+2))^c)}\) for some fixed
\(c\geq 1\). On the kernel this is \(2^{O((k+1)^c)}\).
Including the polynomial preprocessing, this gives the claimed
\(2^{k^{O(1)}}n^{O(1)}\) algorithm for \textsc{Edge Clique Cover}.
\end{proof}

\begin{thebibliography}{1}
\bibitem{gramm}
Jens Gramm, Jiong Guo, Falk H\"uffner, and Rolf Niedermeier.
\newblock Data Reduction and Exact Algorithms for Clique Cover.
\newblock \emph{ACM Journal of Experimental Algorithmics}, 13,
Article 2.2, 2009.
\newblock \href{https://doi.org/10.1145/1412228.1412236}{DOI: 10.1145/1412228.1412236}.
\end{thebibliography}
\end{document}
